\newpage
\section{Các bên liên quan (Stakeholders)}

Việc xác định các bên liên quan giúp hệ thống đáp ứng đúng nhu cầu của từng nhóm người dùng, từ khách hàng vãng lai đến đội ngũ quản trị doanh nghiệp. Các nhóm Stakeholders chính bao gồm:

\begin{center}
\begin{table}[H]
\centering
\renewcommand{\arraystretch}{1.5}
\begin{tabular}{|p{4.5cm}|p{10cm}|}
\hline
\textbf{Nhóm Stakeholder} & \textbf{Vai trò và yêu cầu} \\
\hline
\textbf{Khách hàng (Guest)} & Truy cập website để tìm hiểu năng lực công ty qua trang Giới thiệu, tra cứu giải đáp tại trang Q\&A, và gửi yêu cầu tư vấn/hợp tác thông qua biểu mẫu Liên hệ. \\
\hline
\textbf{Thành viên (Member)} & Khách hàng có tài khoản hệ thống, có thể thực hiện bình luận, đánh giá các bài viết tin tức hoặc sản phẩm, và quản lý thông tin cá nhân. \\
\hline
\textbf{Quản trị viên (Admin)} & \textbf{(Trọng tâm Task 1 \& 2)} Sử dụng dashboard để cấu hình thông tin chung (liên lạc, logo), quản lý danh sách phản hồi từ khách hàng, và cập nhật nội dung bài viết Giới thiệu/Q\&A. \\
\hline
\textbf{Ban lãnh đạo doanh nghiệp} & Theo dõi các báo cáo thống kê về lượt tương tác, phản hồi của khách hàng để đánh giá hiệu quả truyền thông và đưa ra định hướng kinh doanh. \\
\hline
\textbf{Đội ngũ kỹ thuật} & Chịu trách nhiệm bảo trì hệ thống, đảm bảo tính toàn vẹn của cơ sở dữ liệu MySQL và xử lý các vấn đề phát sinh trên nền tảng PHP. \\
\hline
\end{tabular}
\end{table}
\end{center}

\subsection{Phân tích nhu cầu trọng tâm cho Task 1 và Task 2}
Dựa trên bảng phân tích trên, các yêu cầu kỹ thuật đối với vai trò Quản trị viên (Admin) được ưu tiên xử lý bao gồm:
\begin{itemize}
    \item \textbf{Đối với Task 1:} Admin cần một giao diện trực quan để quản lý các mẫu tin liên hệ (Contact submissions), đảm bảo không bỏ lỡ yêu cầu nào từ khách hàng tiềm năng.
    \item \textbf{Đối với Task 2:} Admin cần hệ thống quản trị nội dung linh hoạt để cập nhật thông tin giới thiệu và bộ câu hỏi thường gặp (Q\&A), giúp giảm thiểu thời gian hỗ trợ trực tiếp cho các câu hỏi lặp lại.
\end{itemize}